\mychapter{Trabalhos Relacionados}{cap:rel}\lhead{TRABALHOS RELACIONADOS}

Este capítulo posiciona a dissertação em relação à literatura existente sobre controle de inventário, previsão de demanda intermitente, meta-heurísticas, aprendizado por reforço e seleção adaptativa de políticas. Enquanto o Capítulo~\ref{cap:tec} apresentou os fundamentos conceituais necessários à formulação do problema, este capítulo compara como trabalhos anteriores abordam problemas próximos ao AIPE e em que medida diferem da proposta desta dissertação.

A análise está organizada em cinco eixos: controle de inventário em redes varejistas, previsão de demanda intermitente, métodos de otimização e aprendizado por reforço, seleção adaptativa de políticas e arquiteturas híbridas GA-RL. O objetivo é explicitar o posicionamento da dissertação sem repetir fundamentos já apresentados, destacando sua contribuição como integração metodológica para seleção contextual de políticas de reposição.

\section{Visão Geral dos Métodos e Resultados da Literatura}
\label{sec:visao_geral_literatura}

A literatura de controle de inventário dispõe de uma base consolidada de políticas clássicas, como EOQ, $(s,S)$, estoque-base e jornaleiro~\citep{Axsater2006InventoryControl,Zipkin2000FoundationsInventory}, reconhecidas pela simplicidade, interpretabilidade e facilidade de implantação operacional. Em paralelo, métodos de previsão de demanda intermitente, como Croston~\citep{Croston1972ForecastingIntermittentDemand}, SBA, ARIMA/SARIMAX, LSTM e XGBoost~\citep{Makridakis2020M4Competition,Nasseri2023RetailDemandPrediction}, produzem estimativas que qualificam o comportamento esperado das séries. Contudo, políticas clássicas enfrentam dificuldades em regimes de alta intermitência e demanda heterogênea, e a previsão não determina, por si só, qual política de reposição adotar para cada perfil de série.

A Tabela~\ref{tab:sintese_metodos_literatura} sintetiza as principais famílias de métodos, seus resultados típicos e as limitações que motivam esta dissertação. Na tabela, a coluna ``Família'' agrupa os métodos por linha de pesquisa; ``Exemplos'' lista métodos representativos de cada família; ``Resultado típico'' resume o tipo de ganho ou aplicação mais comum reportado na literatura; e ``Limitação'' indica a lacuna que persiste em relação ao problema tratado nesta dissertação.

\begin{table}[htb]
\centering
\small
\setlength{\tabcolsep}{4pt}
\caption{Síntese dos principais métodos e lacunas da literatura.}
\label{tab:sintese_metodos_literatura}
\begin{tabular}{@{}p{2.8cm}p{3.6cm}p{3.4cm}p{3.4cm}@{}}
\toprule
Família & Exemplos & Resultado típico & Limitação \\
\midrule
Políticas clássicas & EOQ, $(s,S)$, estoque-base, Jornaleiro & Simplicidade, interpretabilidade e facilidade de implantação & Baixa adaptação a séries heterogêneas e regimes intermitentes \\
\addlinespace
Previsão de demanda & Croston, SBA, ARIMA, LSTM, XGBoost & Melhora estimativas e fornece sinais preditivos & Não define qual política de reposição aplicar \\
\addlinespace
Otimização e meta-heurísticas & GA, SA, PSO, DE, MILP & Calibração de parâmetros e redução de custo em cenários específicos & Calibra política ou formulação específica, sem selecionar entre famílias \\
\addlinespace
Aprendizado por reforço & DQN, PPO, SARSA, DRL & Decisões sequenciais em ambientes simulados & Sensível à recompensa, ao ambiente e ao histórico disponível \\
\addlinespace
Seleção e meta-aprendizado & ASP, hiper-heurísticas, AutoML & Escolha de método conforme características da instância & Pouco explorado para seleção de políticas em redes varejistas reais \\
\bottomrule
\end{tabular}
\end{table}

Para tratar essas limitações isoladas, a literatura explorou meta-heurísticas para calibrar parâmetros em espaços não convexos~\citep{Talbi2009Metaheuristics,Yuna2023IntermittentMetaheuristics}, agentes de aprendizado por reforço para aprender decisões sequenciais em ambientes simulados~\citep{Oroojlooyjadid2022BeerGameDQN,Gijsbrechts2022DRLInventoryMSOM} e abordagens de seleção de algoritmos para mapear características de instâncias em métodos adequados~\citep{Rice1976AlgorithmSelection,Kerschke2019AutomatedAlgorithmSelection}. Meta-heurísticas normalmente calibram uma política ou formulação específica, sem resolver automaticamente a escolha entre famílias distintas; agentes de RL demandam históricos longos e são sensíveis à função de recompensa e à modelagem do ambiente; a seleção de algoritmos, por sua vez, tem pouca aplicação direta à escolha de políticas de reposição por série em redes varejistas reais.

A literatura oferece, assim, blocos metodológicos sólidos para previsão de demanda, otimização de parâmetros e controle sequencial de estoque, mas ainda há espaço para integrá-los em um \textit{framework} único que caracterize as séries, avalie empiricamente um portfólio de políticas candidatas e aprenda a selecionar contextualmente a política mais adequada. As seções seguintes detalham como cada frente contribui e onde persistem as lacunas.

\section{Inventário em Redes Varejistas}

A literatura de controle de estoques em redes multi-escalão parte do trabalho fundador de \citet{ClarkScarf1960MultiEchelon}, que estabelece, para estruturas seriadas, que políticas ótimas podem ser decompostas por escalão sob hipóteses específicas. \citet{Axsater2006InventoryControl} e \citet{Zipkin2000FoundationsInventory} consolidam essa base em textos clássicos que formalizam custos de posse e ruptura, políticas de estoque-base (\textit{base-stock}) e medidas de nível de serviço. Em ambientes industriais, modelos Guaranteed Service Model (GSM)~\citep{GravesWillems2000SafetyStock} e suas extensões~\citep{Achkar2024GSMExtensions,Eruguz2016GSMsurvey} constituem a principal família de abordagens para posicionamento estratégico de estoques de segurança em redes gerais, sendo amplamente empregados em planejamento integrado.

O regime de demanda intermitente introduz complexidades que escapam à teoria clássica. \citet{Croston1972ForecastingIntermittentDemand} propôs o primeiro estimador específico para esse regime, modelando separadamente a frequência e o tamanho dos eventos de demanda. \citet{Syntetos2001IntermittentDemand} refinaram essa classificação, introduzindo o diagrama ADI$\times$CV$^2$ que divide os padrões de demanda em quatro quadrantes com recomendações metodológicas distintas, tornando explícita a diferença entre séries de demanda \textit{suave}, \textit{errática}, \textit{intermitente} e \textit{grumosa} (\textit{lumpy}). \citet{DeKok2018Typology} sintetizam o estado da arte em otimização de inventário multi-escalão (MEIO) e identificam como lacunas abertas exatamente as situações de demanda não-estacionária em redes gerais (o contexto desta dissertação).

\section{Previsão de Demanda Intermitente}

A previsão de demanda intermitente apresenta dois desafios estruturais: (i) a frequência de zeros impede o uso de métricas relativas clássicas como MAPE; (ii) a não-estacionariedade invalida hipóteses de modelos como Holt-Winters que pressupõem padrões suaves. Evidência empírica em competições de previsão (M4/M5) indica que modelos baseados em árvores de decisão e gradiente \textit{boosting} frequentemente têm desempenho melhor que redes LSTM em séries com alta intermitência, especialmente quando o volume de dados históricos por SKU é limitado~\citep{Nasseri2023RetailDemandPrediction}.

No plano das métricas, \citet{Hyndman2006MASE} introduziram o MASE como substituto livre de escala para o MAPE: é definido com demanda real zero, é comparável entre séries de escalas distintas e não sofre da assimetria do MAPE. \citet{Makridakis1993SMAPE} propôs o sMAPE como variante simétrica; a competição M4 adotou o MASE como métrica primária, consolidando seu uso \citep{Makridakis2020M4Competition}.

\citet{TadayonradNdiaye2023KPIDemandForecasting} argumentam que essas métricas, embora mais adequadas que o MAPE em séries intermitentes, não capturam o impacto sobre custos operacionais e propõem KPIs alinhados ao custo e ao nível de serviço. Esta dissertação segue essa direção: o módulo de previsão alimenta características ao AIPE em vez de ser avaliado isoladamente.

A previsão de demanda, entretanto, não define sozinha a política de reposição. Ela fornece sinais auxiliares sobre o comportamento esperado da série, mas a decisão de quando pedir e quanto pedir depende de uma política de controle, cujo ajuste ou aprendizado por meta-heurísticas e por aprendizado por reforço é discutido na seção seguinte.

\section{Meta-Heurísticas e Aprendizado por Reforço em Políticas de Reposição}

Enquanto a seção anterior discutiu métodos de previsão de demanda, esta seção trata de trabalhos que aplicam meta-heurísticas e aprendizado por reforço ao controle de inventário e à parametrização de políticas de reposição. Nesse contexto, o objetivo não é prever diretamente a demanda, mas decidir quando emitir pedidos e em que quantidade, ou calibrar parâmetros de políticas como $(s,S)$, EOQ ou regras equivalentes. Essa distinção é importante para esta dissertação, pois a previsão é usada como sinal auxiliar, enquanto a decisão central do AIPE é selecionar políticas de reposição.

A otimização de parâmetros de políticas de inventário por meta-heurísticas tem sido extensivamente estudada. \citet{Talbi2009Metaheuristics} consolida a base teórica de GA, SA, PSO e outros métodos de busca para problemas combinatórios e de otimização em espaços de parâmetros.

No contexto de demanda intermitente, \citet{Yuna2023IntermittentMetaheuristics} compararam SA e Busca Tabu na otimização de parâmetros de uma política $(s,S)$ em sete conjuntos com CV$>$1.5; \citet{Erkayman2025MarkovInventory} modelaram demanda como cadeia de Markov e otimizaram parâmetros com GA e PSO. Esses estudos mostram que a calibração automática pode melhorar o desempenho de uma política específica, isto é, resolvem \textit{como} calibrar uma política já escolhida, mas não \textit{qual} política adotar dado o perfil da série. Além disso, restringem-se a modelos de único escalão sem comparação entre famílias distintas de política.

Na frente de aprendizado por reforço, os trabalhos relacionados tratam a reposição como um problema sequencial de decisão, não como previsão. \citet{Oroojlooyjadid2022BeerGameDQN} demonstraram que DQN aprende políticas de reposição no \textit{Beer Game} multi-escalão sem especificar a distribuição de demanda, mas assumem séries longas e demanda contínua. \citet{Geevers2024MEIODRL} aplicaram PPO a três topologias de rede com desempenho melhor que o \textit{benchmark}, porém em ambientes sintéticos com CV~$<1$. \citet{Liu2024HAPPO} introduziram HAPPO para reduzir simultaneamente custo e efeito \textit{bullwhip} em ambientes descentralizados, mas não abordam demanda intermitente.

A limitação transversal desses trabalhos está na distância entre calibrar ou treinar uma política específica e selecionar contextualmente entre famílias distintas de políticas: meta-heurísticas calibram bem uma formulação já escolhida, e agentes de RL podem aprender boas decisões em um ambiente específico, mas nenhuma dessas abordagens responde diretamente \textit{qual} política deve ser adotada para uma série loja-produto com características operacionais específicas. Soma-se a isso uma limitação prática do RL: a exigência da ordem de $10^6$ interações agente-ambiente, isto é, passos sequenciais em que o agente observa o estado do estoque, escolhe uma ação de reposição, recebe a recompensa associada ao custo e ao nível de serviço e atualiza sua política~\citep{Oroojlooyjadid2022BeerGameDQN}. Essa escala é incompatível com séries históricas curtas de varejo (38 ciclos no \textit{dataset} desta dissertação) e é a razão estrutural pela qual RL isolado degenera no regime estudado. Essa lacuna justifica tratar meta-heurísticas e agentes de RL como componentes de um portfólio de políticas candidatas, e não como solução única para todas as séries.

\section{Seleção de Políticas, Hiper-Heurísticas e Meta-Aprendizado}
\label{sec:aas_literatura}

O problema de determinar qual algoritmo usar dada uma instância foi formalizado por \citet{Rice1976AlgorithmSelection} como o \textit{Algorithm Selection Problem} (ASP). Dado um espaço de instâncias $\mathcal{I}$, um espaço de algoritmos $\mathcal{A}$ e uma medida de desempenho $f$, o objetivo é encontrar $s^*: \mathcal{I} \rightarrow \mathcal{A}$ tal que $s^*(i) = \arg\max_{a \in \mathcal{A}} f(a, i)$. O AIPE é uma instância direta do ASP aplicada ao controle de inventário: as instâncias são séries de demanda, os algoritmos são as políticas candidatas e $f$ é o CTI minimizado sujeito a NS mínimo.

O teorema \textit{No Free Lunch} (NFL) de \citet{Wolpert1997NoFreeLunch} fornece a justificativa teórica: nenhum algoritmo domina todos os outros na média sobre todas as instâncias. Isso implica que qualquer ganho de desempenho consistente exige explorar estrutura específica da classe de problemas. As características operacionais usadas como entrada pelo AIPE (ADI, CV$^2$, \textit{burstiness}, entropia, sazonalidade) são exatamente essa estrutura discriminativa.

A implementação prática do ASP segue o paradigma \textit{características de instância $\rightarrow$ modelo supervisionado $\rightarrow$ seleção de algoritmo}. \citet{Kerschke2019AutomatedAlgorithmSelection} consolidam esse paradigma em \textit{survey} abrangente; \citet{Kerschke2019ELASelection} demonstram empiricamente que características de paisagem combinadas com classificadores supervisionados selecionam automaticamente o otimizador mais adequado por instância, a mesma arquitetura conceitual adotada pelo AIPE ao usar características operacionais de demanda como entrada. \citet{Kotthoff2014AlgorithmSelection} mostra que seletores por instância podem obter ganhos frente a políticas fixas em problemas combinatórios.

No campo das \textit{hyper-heuristics}, \citet{Drake2020SelectionHyperHeuristics} e \citet{Dokeroglu2024HyperHeuristics} consolidam o estado da arte: \textit{selection hyper-heuristics} escolhem, a cada estado do problema, qual heurística aplicar. O AIPE é análogo operando no nível de série, mas com escopo distinto: seleciona \textit{offline} a política a partir do histórico e do perfil de demanda da combinação loja-produto, antes da execução operacional. Uma seleção \textit{online} por estado de estoque implicaria trocar heurísticas a cada ciclo ou decisão de reposição, exigindo um controlador em tempo real; nesta dissertação, essa responsabilidade é substituída por recomendações periódicas de política por série.

Em inventário especificamente, \citet{Preil2022BanditInventory} modelam a seleção como \textit{multi-armed bandit}, e \citet{ContextualRL2024SupplyChain} adaptam políticas de RL via informação de contexto. O AIPE se alinha a essa direção ao deslocar o foco da adaptação de parâmetros de uma política fixada para a seleção contextual entre famílias de política. \citet{Qi2023EndToEndInventory} mostram em \textit{Management Science} que características operacionais podem sustentar aprendizado de políticas de inventário em escala industrial.

\smallskip
A literatura de ASP, \textit{hyper-heuristics} e meta-aprendizado converge para o mesmo ponto: a política mais adequada depende da instância, essa dependência pode ser aprendida e, em contextos apropriados, um meta-modelo treinado com evidência empírica pode obter ganhos frente a políticas fixas. Neste trabalho, esse princípio é aplicado ao problema de inventário varejista sob demanda intermitente.

\section{Políticas Híbridas GA-RL}

A frente mais próxima combina busca global de meta-heurísticas com adaptação local do RL. \citet{Zabraoui2025GaDQN} é o artigo-base que inspirou a inclusão da arquitetura GA-DQN no portfólio desta dissertação. Os autores introduziram GA-DQN sobre o conjunto Walmart~M5 (CV típico $<1$): o GA otimiza parâmetros globalmente enquanto o DQN refina decisões a cada ciclo, elevando o NS de 61\% para 94\% com redução de custo. Esta dissertação incorpora GA-DQN como \textit{uma das doze políticas avaliadas} no portfólio, mas posiciona o \textit{Policy Selection Engine} como contribuição principal: o objetivo não é provar que GA-DQN vence em todos os regimes, mas aprender \textit{quando} cada política do portfólio apresenta vantagem relativa. Esse enquadramento é compatível com a evidência empírica observada, pois nenhuma política domina todas as métricas em todos os regimes. Também preserva abertura metodológica: o PSE pode incorporar novas políticas ao \textit{benchmark} desde que seus desempenhos sejam avaliados e os rótulos do meta-modelo sejam atualizados.

A variante GA-PPO, que substitui o DQN por uma política contínua baseada em PPO~\citep{Vanvuchelen2020PPOJointReplenishment,Gijsbrechts2022DRLInventoryMSOM}, segue a mesma motivação de combinar busca evolutiva global com refinamento por gradiente de política, mas apresenta dinâmicas distintas em regimes de alta volatilidade. Ambas as arquiteturas híbridas figuram como políticas especialistas dentro do portfólio de doze candidatas avaliadas empiricamente nesta dissertação. A presença de GA-DQN e GA-PPO no portfólio possibilita ao PSE identificar, para cada perfil operacional de demanda, se uma abordagem híbrida oferece melhor compromisso entre custo e nível de serviço ou se políticas clássicas permanecem mais adequadas naquele recorte.

\section{Posicionamento da Dissertação}

A Tabela~\ref{tab:relwork_intro} posiciona esta dissertação em relação aos trabalhos mais próximos segundo seis dimensões. A coluna ``Trabalho'' identifica o estudo analisado; ``Base usada'' diferencia avaliações conduzidas em ambientes sintéticos, conjuntos de teste e dados transacionais reais; ``CV ou regime'' resume o nível de variabilidade ou o regime de demanda considerado; ``Políticas avaliadas'' informa quantas alternativas de reposição ou controle foram comparadas; ``Multinível'' indica se o estudo considera mais de um escalão da cadeia de suprimentos; e ``Indicadores avaliados'' informa o número de métricas operacionais utilizadas. Essa comparação permite situar a dissertação quanto à escala experimental, ao tipo de dado e à abrangência metodológica.

\begin{table}[htb]
\centering
\footnotesize
\setlength{\tabcolsep}{4pt}
\caption{Posicionamento desta dissertação em relação aos trabalhos relacionados.}
\label{tab:relwork_intro}
\begin{tabular}{@{}p{3.1cm}p{2.1cm}cccc@{}}
\toprule
Trabalho &
\makecell[c]{Base\\usada} &
\makecell[c]{CV ou\\regime} &
\makecell[c]{Políticas\\avaliadas} &
\makecell[c]{Multi-\\nível} &
\makecell[c]{Indicadores\\avaliados} \\
\midrule
\footnotesize\citealt{Oroojlooyjadid2022BeerGameDQN} & Sintética & $<1$ & 3 & Sim & 2 \\
\citealt{Vanvuchelen2020PPOJointReplenishment} & Sintética & $<1$ & 4 & Não & 1 \\
\citealt{Gijsbrechts2022DRLInventoryMSOM} & Sintética & $<1$ & 6 & Sim & 2 \\
\citealt{Yuna2023IntermittentMetaheuristics} & Testes intermitentes & $>1.5$ & 2 & Não & 2 \\
\citealt{Erkayman2025MarkovInventory} & Sint. Markov & $>1.5$ & 2 & Não & 1 \\
\citealt{Liu2024HAPPO} & Sintética & $<1$ & 4 & Sim & 2 \\
\citealt{Zabraoui2025GaDQN} & Real (Walmart M5) & $<1$ & 7 & Sim & 3 \\
\midrule
Esta dissertação (Experimento 1) & Real (PB) & $\mathbf{2.5{-}5.4}$ & 12 & Sim & 5 \\
Esta dissertação (Experimento 2) & Real (BA) & Lumpy; etapa nacional futura & 12 & Sim & 5 + testes \\
\bottomrule
\end{tabular}
\end{table}

A revisão da literatura aponta três lacunas convergentes. A primeira é empírica: embora trabalhos recentes já incluam bases reais de varejo, como o Walmart M5 em \citet{Zabraoui2025GaDQN}, ainda há poucos estudos que avaliam políticas de reposição em dados reais de varejo brasileiro com demanda altamente intermitente e séries loja-produto heterogêneas. A segunda é metodológica: trabalhos existentes tendem a otimizar ou aprender políticas específicas, mas não formulam a escolha da política como um problema de seleção contextual por perfil operacional da série. A terceira é operacional: ainda é pouco explorada a integração, em um mesmo \textit{pipeline} reprodutível, entre caracterização das séries, simulação comparável de políticas, validação estatística, geração de rótulos de política dominante e treinamento de um seletor supervisionado. Essas lacunas motivam a formulação do AIPE como abordagem integrada de avaliação e recomendação de políticas.

As seções anteriores indicam que decisões de reposição em redes varejistas heterogêneas dependem do comportamento operacional de cada série loja-produto. Políticas clássicas, políticas parametrizadas por meta-heurísticas e políticas aprendidas por reforço oferecem alternativas relevantes, mas seu desempenho tende a variar conforme frequência, variabilidade e volume da demanda. Do mesmo modo, a previsão de demanda contribui como sinal auxiliar, mas não resolve, isoladamente, a escolha da política de reposição mais adequada.

Assim, a contribuição desta dissertação não está em propor novos algoritmos de otimização, aprendizado por reforço ou reposição: GA, SA, PSO, DE, DQN, PPO, SARSA e as arquiteturas híbridas GA-RL já existem na literatura e são utilizados como componentes conhecidos de um portfólio experimental, avaliados empiricamente sob o mesmo protocolo. O papel do AIPE é organizar esses componentes em um \textit{framework} metodológico de seleção contextual, no qual séries loja-produto são caracterizadas, políticas candidatas são avaliadas sob um ambiente comum de simulação e os resultados são usados para recomendar políticas de acordo com o perfil operacional observado, sobre dados reais de uma rede varejista brasileira.

Nesse sentido, a comparação entre múltiplas políticas é necessária para construir o \textit{benchmark} e gerar rótulos de política dominante, mas não constitui, isoladamente, a contribuição central. O ponto central é avaliar se a seleção contextual, aplicada a esses dados reais, reduz o custo de inventário em relação à adoção de uma política única global, preservando critérios mínimos de serviço e analisando estabilidade operacional. O treinamento do \textit{Policy Selection Engine} a partir dessa evidência completa essa contribuição.

A partir desse posicionamento, o AIPE é tratado como um \textit{framework} metodológico de seleção contextual de políticas, e não como uma nova política universal de reposição. O próximo capítulo apresenta a metodologia adotada para operacionalizar essa proposta: avaliação das políticas candidatas em simulação comum, comparação por métricas de custo, serviço e estabilidade, geração de rótulos de política dominante e treinamento do \textit{Policy Selection Engine} (PSE).

